\section{Preparation precision at full endpoints}
\label{sec:precision}

An exactly full coordinate can transfer immediately, whereas a strict
deficit may expose a rate barrier. We quantify this distinction for
$n\geq1$ modules with positive proportional coefficients. For a maintained
state $p$, suppose a request may reveal any state $z$ satisfying
$(p_i-\epsilon)_+\leq z_i\leq p_i$. Here $\epsilon>0$ is an absolute
preparation deficit in the common work units. This is a bounded perturbation
at the request, not a stream of additional normal-operation disturbances.
Monotonicity makes $p^\epsilon=((p_i-\epsilon)_+)_i$ its worst corner.
Define
\begin{equation}
 F_-(p)=\lim_{\epsilon\downarrow0}F(p^\epsilon),\quad
 C_-(H)=\inf_{F_-(p)\leq H}L(p),\quad
 C_\epsilon(H)=\inf_{F(p^\epsilon)\leq H}L(p),
 \label{eq:precision-costs}
\end{equation}
where $L(p)=\sum_i\gamma_i p_i$ and empty infima are infinite. The limit
exists by monotonicity. Maintenance is charged to the nominal state $p$.

\begin{proposition}[One-sided exit envelope]
\label{prop:precision-oracle}
Let $\mathcal P$ contain the permutations satisfying
$b(S_{k-1})>d_{\pi_k}$ at every stage, where
$S_{k-1}=\{\pi_1,\ldots,\pi_{k-1}\}$ and $d_i=\gamma_iM_i$.
For each such order, define $G_\pi(p)=t_n$ by
\begin{equation}
 t_0=0,\qquad
 t_k=\frac1{\gamma_{\pi_k}}
 \log\frac{b(S_{k-1})e^{\gamma_{\pi_k}t_{k-1}}
                  -\gamma_{\pi_k}p_{\pi_k}}
                 {b(S_{k-1})-d_{\pi_k}}.
 \label{eq:precision-recurrence}
\end{equation}
Then $F_-(p)=\min_{\pi\in\mathcal P}G_\pi(p)$.
If cold exit is finite, $F_-$ is finite and continuous on the closed
preparation box. Otherwise it is identically infinite.
Always $F\leq F_-$, with equality at every subfull state.
\end{proposition}

\begin{proof}
Every $p^\epsilon$ is subfull. Theorem~\ref{thm:seriality} and the
strict finite-time rate barrier therefore allow exactly the orders in
$\mathcal P$, independently of $\epsilon$. Equation~\eqref{eq:proportional-stage}
gives \eqref{eq:precision-recurrence}. Each recurrence is continuous on the
closed box, with positive denominators and nonnegative stage durations.
Consequently its finite minimum commutes with the limit. The same orders
are feasible from cold, so the family is nonempty exactly when cold exit
is finite. Monotonicity gives $F\leq F_-$, and the identical order formula
applies directly to $F$ at subfull states.
\end{proof}

The strict inequalities include initially full coordinates. They can have
zero-length stages at time zero when accessible; a full coordinate left
waiting decays and requires preparation again. If the full coordinates
admit a strictly accessible order from the empty set, placing them first
at zero time and following an optimal original schedule proves
$F_-(p)=F(p)$. In particular this holds when every full coordinate has
$d_i<s$.

\begin{proposition}[Maintenance below the normal budget]
\label{prop:precision-slack}
If $C(H)<s$, then $C_-(H)=C(H)$.
\end{proposition}

\begin{proof}
At an original minimum with upkeep below $s$, each full coordinate has
$d_i<s$. The preceding zero-length prefix construction makes that state
feasible for $C_-$, giving equality since $F\leq F_-$ implies $C\leq C_-$.
\end{proof}

At $C(H)=s$, failure of equality can occur only when every original
minimum is a full singleton of cost $d_i=s$. Any other minimum has all
full losses strictly below $s$ and supplies the same zero-length prefix.
The usual total-storage bound and maintenance construction apply to the
upward closed limiting ready set, giving optional rate $s-C_-$ when
$C_-\leq s$ and excluding indefinite readiness otherwise.

\begin{proposition}[Concentration with preparation tolerance]
\label{prop:precision-concentration}
Whenever the respective finite-deadline problem is feasible, every
minimizing state for $C_-(H)$ and every minimizing nominal state for
$C_\epsilon(H)$ has at most one partial coordinate. This holds for unequal
positive coefficients and at every source budget.
\end{proposition}

\begin{proof}
For the limiting problem use $q=p$ with caps $U_i=M_i$. For fixed positive
tolerance, an optimizer cannot have $0<p_i\leq\min(\epsilon,M_i)$:
replacing it by zero preserves its worst corner and lowers upkeep. Hence
use $q_i=(p_i-\epsilon)_+$ with caps $U_i=(M_i-\epsilon)_+$ and upkeep
$\sum_i\gamma_iq_i+\sum_{i:q_i>0}\gamma_i\epsilon$.
Zero q corresponds to zero p, while $q_i=U_i>0$ corresponds to $p_i=M_i$.
Both problems therefore use the strict-order oracle with capped q; the
additional positive-support charges stay constant under interior variations.
Their minima exist when feasible: the limiting feasible set is compact
by Proposition~\ref{prop:precision-oracle}, and the fixed-tolerance
nominal constraint is continuous by applying that oracle to its subfull
worst corner.

For fixed preparation q, a stage map $\phi(t)$ from
\eqref{eq:precision-recurrence} satisfies
\[
 \phi'\geq1,\qquad
 \phi''=-\gamma\phi'(\phi'-1)\leq0,\qquad \phi(0)\geq0.
\]
The composition G of any intervening stages is increasing and concave,
with $G(0)\geq0$. Its tangent inequality gives $G(t)\geq tG'(t)$.

Suppose two coordinates j,k are interior to their positive caps in an
attaining strict order. Keep the start of j fixed, vary its completion t,
and keep the completion of k fixed. All intervening preparations remain
fixed, with G mapping t to the start of k. Writing
$\alpha=\gamma_j$, $\beta=\gamma_k$, $D_j=b_j-d_j>0$, the pair upkeep
up to constant terms is
\[
 c(t)=-D_j e^{\alpha t}+b_k e^{\beta G(t)}.
\]
An open interval preserves both interior caps and positive stage durations;
the same completion of k reproduces the full suffix and its deadline.
At a stationary point,
$\alpha D_j e^{\alpha t}=\beta b_k G'e^{\beta G}$.
If $\beta G'\geq\alpha$, then $\beta G\geq\alpha t$, contradicting
this equality because $b_k\geq b_j>D_j$. Thus $\beta G'<\alpha$, and
\[
 c''(t)=\alpha D_j e^{\alpha t}
       \left(\beta G'+\frac{G''}{G'}-\alpha\right)<0.
\]
Every stationary point is a strict local maximum. The proposed minimum
cannot have two interior coordinates. Under the nominal conversion above,
this is exactly the asserted concentration.
\end{proof}

The worst-corner state can itself have several physically partial
coordinates: nominally full modules have preparation $M_i-\epsilon$ at
the request. The result concerns the nominal design, and does not remove
their post-request preparation or justify treating them as free transfers.

\begin{proposition}[Positive tolerance and its limit]
\label{prop:precision-limit}
If cold exit is finite and $H>0$, then
$\lim_{\epsilon\downarrow0}C_\epsilon(H)=C_-(H)$.
For every positive $\epsilon$, whenever $C(H)>0$,
\begin{equation}
 C_\epsilon(H)\geq C(H)+\gamma_{\min}\epsilon,
 \qquad \gamma_{\min}=\min_i\gamma_i>0.
 \label{eq:precision-margin}
\end{equation}
\end{proposition}

\begin{proof}
Monotonicity gives $F_-(p)\leq F(p^\epsilon)$, hence $C_\epsilon\geq C_-$.
Choose a minimizing state for $C_-$. If its time is strictly below $H$,
it is feasible at sufficiently small positive tolerance. Otherwise its
time equals $H>0$, so it is not all-full. Moving it an arbitrarily small
distance toward the all-full vector strictly lowers an attaining order's
time: increasing a coordinate strictly advances its own stage, and every
later recurrence is strictly increasing in its preceding time. The new
state has time slack and arbitrarily close upkeep. It is feasible for all
sufficiently small tolerances, proving the opposite limit inequality.

For \eqref{eq:precision-margin}, any feasible nominal state has
$q=p^\epsilon$ with $L(q)\geq C(H)>0$. Some $q_i>0$, so
$p_i=q_i+\epsilon$, while every other $p_j\geq q_j$. Thus
$L(p)\geq L(q)+\gamma_{\min}\epsilon$; take the infimum.
\end{proof}

The condition $H>0$ is essential: positive deficits make a zero deadline
infeasible even though the all-full state has $F_-=0$ when cold exit is
finite. For positive deadlines, strict optional slack survives sufficiently
small positive tolerance, although its exact optimal rate can decrease.
Every originally critical instance $C=s>0$ loses indefinite readiness at
any positive tolerance by \eqref{eq:precision-margin} and the storage
bound. That elementary margin obstruction is general; the next example's
distinctive feature is a nonvanishing limiting cost jump.

\begin{example}[A critical endpoint jump with finite cold exit]
\label{ex:precision-jump}
Take
\[
 s=\gamma_A=\gamma_B=1,\quad
 M_A=\tfrac12,\quad M_B=1,\quad a_A=a_B=1,
 \qquad H=\log(4/3).
\]
Only the order A,B is strictly accessible. Its recurrence is
$\exp F_-(p)=4-4p_A-p_B$, so cold exit takes $\log4$.
The original state $(0,1)$ transfers B immediately and finishes A at
capacity two in time $\log(4/3)$, whereas $F_-(0,1)=\log3$.
The limiting deadline requires $4p_A+p_B\geq8/3$, implying
$p_A+p_B\geq8/3-3p_A\geq7/6$. Equality holds only at
$(1/2,2/3)$. Every original ready state with $p_B<1$ obeys this same
constraint, while B-full states cost at least one. Consequently
\[
 C(H)=1=s,\qquad C_-(H)=\tfrac76>s.
\]
For fixed tolerance, both coordinates of $q=p^\epsilon$ must be positive
to satisfy $4q_A+q_B\geq8/3$. Nominal upkeep is then
$q_A+q_B+2\epsilon$, minimized by
$q_A=1/2-\epsilon$, $q_B=2/3+4\epsilon$. The latter obeys
$q_B\leq1-\epsilon$ exactly when $\epsilon\leq1/15$. Hence
\[
 C_\epsilon(H)=
 \begin{cases}
  7/6+5\epsilon,&0<\epsilon\leq1/15,\\
  +\infty,&\epsilon>1/15.
 \end{cases}
\]
For larger tolerance even the all-full nominal state is infeasible.
The limiting increase $1/6$ therefore persists as precision improves.
Two modules are necessary under finite cold exit: a single module requires
$s>d$ and has the continuous formula
$F(p)=\gamma^{-1}\log[(s-\gamma p)/(s-d)]$ through its full endpoint.
\end{example}
